\bigskip

% course 9, page 1
% \section*{Curs 9 - Pagina 1}


\begin{center}
\fbox{\textbf{ALGEBRA}}\\[2pt]
\fbox{\textbf{CURS 9}}
\end{center}
\begin{flushright}
\fbox{\textit{29-XI-1994}}
\end{flushright}



\bigskip
\noindent\underline{Subgr.\ generat de o mulţ.\ de elem.\ ale unui grup}

\noindent\underline{Grupul comutator (derivat) al unui grup}

\noindent\underline{Grupul diedral $\Delta_n$, $n\geq 3$}

\bigskip
\bigskip
\noindent Fie $(G,\cdot,e)$ grup, $M\subseteq G$, $M\neq\varnothing$

\noindent\underline{Problemă}: \ Să se descrie c.m.\ mic subgr.\ al lui $G$

\noindent care conţ.\ pe $M$.

\bigskip
\noindent\underline{Prop.} \ Dacă $H_i$, $i\in I$ fam.\ de subgr.\ ale lui $G$ şi

\[
H=\bigcap_{i\in I} H_i \quad \text{atunci} \quad H\leq G.
\]

\bigskip
\noindent\underline{Dem} \ BANALĂ!

\noindent $H\neq\varnothing$ \quad $e\in H_i$ $(\forall)\,i \Rightarrow e\in H$

\noindent Dacă $x,y\in H$ şi $\left.\begin{aligned}x,y&\in H_i\ (\forall)\,i \Rightarrow xy\in H_i\ (\forall)\,i\\ x^{-1}&\in H_i\ (\forall)\,i\end{aligned}\right\} \Rightarrow xy,x^{-1}\in H$

\bigskip
\noindent\underline{Obs.} \ (1) \ Dacă $H\leq G$ şi $M\subseteq H \Rightarrow \langle M\rangle \leq H$ \ unde

\[
\langle M\rangle = \{x_1^{\varepsilon_1}x_2^{\varepsilon_2}\cdots x_m^{\varepsilon_m} \mid m\geq 1,\ x_i\in M,\ \varepsilon_i=\pm 1\ (\forall)\,i\}
\]

\noindent (2) $\langle M\rangle \leq G$

\noindent (3) $M\subseteq \langle M\rangle$ \quad (elem.\ lui $M$ = generatori; $\langle M\rangle$ = subgr.\ lui $G$ generat de $M$)

\noindent (4) $\displaystyle\bigcap_{M\subseteq H\subseteq G} H$ = c.m.\ mic subgr.\ al lui $G$ care conţ.\ pe $M$ \quad (de dem.)

\bigskip
\noindent $\langle\varnothing\rangle := \langle\{e\}\rangle$

\bigskip
\noindent Dacă $M=\{a_1,\ldots,a_n\}$ atunci $\langle\{a_1,\ldots,a_n\}\rangle = \langle a_1,\ldots,a_n\rangle$

\bigskip
\noindent\underline{Def.} \ Fie $G$ grup.\ Spunem că $G$ este \underline{finit generat}

\noindent dacă $\exists\,a_1,\ldots,a_m\in G$ a.î.\ $G=\langle a_1,\ldots,a_m\rangle$

\bigskip
\noindent $G$ \underline{ciclic} dacă $\exists\,a\in G$ a.î.\ $\langle a\rangle=G$

\bigskip
\noindent $S_n = \langle \tau \mid \tau\text{-transp.}\rangle$

\noindent Dacă $M=\{\tau\in S_n \mid \tau \text{ transp.}\}$

\bigskip
\noindent\underline{Ex}/ \quad $\mathbb{Z} = \langle 1\rangle$

\bigskip
\noindent\underline{Ex}/ \quad Dem.\ că $\big\langle (1,2),(1,3),\ldots,(1,n)\big\rangle = S_n$



% course 9, page 2
\section*{Curs 9 - Pagina 2}

\noindent\underline{\textbf{Grupul comutator (derivat) al unui grup.}}

\bigskip
\noindent Fie $(G,\cdot,e)$ grup.

\noindent $a,b\in G$ \quad $[a,b] \overset{\text{def}}{=} aba^{-1}b^{-1}\in G$

\noindent s.n.\ \underline{comutator} (comutatorul

\noindent perechii de elem.\ $a,b$)

\bigskip
\noindent $[a,b]\,ba = ab$

\noindent $c\in G \Rightarrow ab=cba$

\bigskip
\noindent\fbox{$G' = \big\langle aba^{-1}b^{-1} \mid a,b\in G\big\rangle$} \quad (mulţ.\ tuturor comutat.)

\bigskip
\noindent Se obs.\ că $[a,b]^{-1} = (aba^{-1}b^{-1})^{-1} = bab^{-1}a^{-1} = [b,a]$

\noindent Deci $G' = \{$prod.\ finite de comutatori$\}$

\bigskip
\noindent Dacă $a,b,c\in G \Rightarrow a[b,c]a^{-1} = abcb^{-1}c^{-1}a^{-1} =$

\[
= aba^{-1}\cdot aca^{-1}\cdot ab^{-1}a^{-1}\cdot ac^{-1}a^{-1} = \big[aba^{-1},\,aca^{-1}\big]
\]

\noindent $\Rightarrow G'\trianglelefteq G$

\bigskip
\noindent Arăt că $axa^{-1}\in G'$ $(\forall)\,a\in G$, $(\forall)\,x\in G'$

\[
x=[b_1,c_1]\cdots[b_n,c_n]
\]
\[
\Rightarrow axa^{-1} = \big[ab_1a^{-1},ac_1a^{-1}\big]\cdots\big[ab_na^{-1},ac_na^{-1}\big]
\]

\noindent (se intercalează produse de forma $a^{-1}a$)

\bigskip
\noindent\underline{Obs}/ \quad \fbox{$\dfrac{G}{G'}$ grup abelian}

\bigskip
\noindent Dacă $\hat a,\hat b\in \dfrac{G}{G'}$ arăt că $\hat a\hat b=\hat b\hat a$

\[
\Leftrightarrow ab\equiv_{d} ba \pmod{G'} \ \Leftrightarrow \ ab(ba)^{-1}\in G'
\]
\[
aba^{-1}b^{-1}\in G' \ \Leftrightarrow \ [a,b]\in G' \quad (\forall)\text{ comut.}\in G'
\]

\noindent adev.\ $\Rightarrow \hat a\hat b=\hat b\hat a$

\noindent $\Rightarrow \dfrac{G}{G'}$ abelian

\bigskip
\noindent\underline{Prop.}: \ \fbox{Dacă $N\trianglelefteq G$ a.î.\ $\dfrac{G}{N}$ abelian $\Rightarrow G'\subseteq N$}\footnote{în original prima literă a acestei propoziţii pare ,,$H$'' în loc de ,,$N$''; corectat aici la $N$ pentru consecvenţă cu restul demonstraţiei, care foloseşte exclusiv $N$.}

\bigskip
\noindent\underline{Dem} \ Fie $a,b\in G$. Atunci $\hat a\hat b=\hat b\hat a$ unde $\hat a=aN$

\noindent $\hat b=bN$

\[
\Leftrightarrow ab\equiv_{d} ba \pmod{N} \ \Leftrightarrow \ ab(ba)^{-1}\in N
\]
\[
[a,b]\in N \ \Rightarrow \ (\forall)\text{ comut.}\in N \ \Rightarrow \ G'\subseteq N
\]

\bigskip
\noindent $\dfrac{G}{G'} = G^{ab}$ \ abelianizatul lui $G$.
